\documentstyle[% 


righttag% 


,11pt% 


,amssymb% 


,russian%               remove in Eng. 


,izvru%                 Set izven% in Eng. 


]{amsart} 


 


% seek 0,5


\newcommand{\Grad}{\operatorname{grad}} 


\newcommand{\Div}{\operatorname{div}} 


\newcommand{\Rot}{\operatorname{rot}} 


%\newcommand{\arccos}{\operatorname{arccos}} 


\newcommand{\const}{\operatorname{const}} 


\newcommand{\tg}{\operatorname{tg}} 


\newcommand{\tts}[1]{} 


 


\theoremstyle{plain} 


\theorembodyfont{\it} 


\newtheorem{thmnonum}{\hskip\parindent Теорема} 


\renewcommand{\thethmnonum}{} 


 


%\hoffset=-15mm%                  For Eng. 


\setcounter{page}{24} 


\god{2004} 


\nomer{2~(501)} %                        Remove in Eng. 


%% remove ^^^^^^ 


%\nomer{Vol.\,48, No.\,2}%               Set in Eng. 


%\str{\pageref{begin}--\pageref{end}}%  Set in Eng. 


\udk{517.544} 


 


\author{\it И.Т.\,Денисюк} 


% NB:   ^^ remove in Eng. 


\title{Задача сопряжения решений уравнения Ламе \\


в областях с кусочно-гладкими границами} 


 


 


\date{} 


%\pagestyle{myheadings}%         Set in Eng. 


\pagestyle{plain} %              Remove in Eng. 


\begin{document} 


%\thispagestyle{plain}%          Set in Eng. 


\maketitle 


\label{begin} 


 


 


 


 


Решение двумерных задач сопряжения аналитических функций в заданных и  


аффинно-пре\-образованных областях с кусочно-гладкими границами [1], [2]  


позволили получить решение ряда прикладных задач механики сплошной  


среды [3]--[9]. В данной статье изучается задача сопряжения решений уравнений  


Ламе [10] в трехмерных областях с кусочно-гладкими границами. 


\medskip 


 


{\bf Постановка задачи.} Пусть в трехмерном пространстве $R_3$ содержится  


односвязная конечная область $V_1$, ограниченная кусочно-гладкой поверхностью  


$S$, которая содержит непересекающиеся гладкие замкнутые особые линии  


(множества угловых точек) и не принадлежащие им конические точки. 


 


Построим векторные функции $\ov u_i(M)=\{u_{1i}(M),u_{2i}(M),u_{3i}(M)\}$ 


($i=\ov{0,1}$, $M=M(x_1,x_2,x_3)$), 


определенные в соответствующих областях $V_i$, где 


$V_0=R_3\setminus V_1$, удовлетворяющие уравнению Ламе [10] 


\begin{equation} 


(\lambda_i+2\mu_i)\Grad\Div\ov u_i(M)-\mu_i\Rot\ov u_i(M)=0, 


\end{equation} 


где $\lambda_i=\frac{2\nu_i\mu_i}{1-2\nu_i}$, $0<\nu_i<0.5$; 


$\mu_i\in(0,\infty)$, $\mu_i\ne\infty$, \\ 


и условиям сопряжения на поверхности $S$ раздела областей $V_0$ и $V_1$: \\ 


в точках гладкости 


\begin{equation} 


\ov v{}^-_0(M)-\ov u{}^+_1(M)=0,\quad 


N^-_{n0}[\ov u_0(M)]-N^+_{n1}[\ov u_1(M)]=0, 


\end{equation} 


в особых точках $M^0$ 


\begin{equation} 


\lim_{M\to M^0}(\ov u{}^-_0(M)-\ov u{}^+_1(M))=0,\quad 


\lim_{M\to M^0}(N^-_{n0}[\ov u_0(M)]- 


N^+_{n1}[\ov u_1(M)])=0, 


\end{equation} 


где оператор $N_{ni}[\ ]$ действует согласно правилу 


$$ 


N_{ni}[\ov u_i(M)]=2\mu_i\frac{\partial\ov u_i(M)}{\partial n} 


+\lambda_i\ov n\Div\ov u_i(M)+\mu_i[\ov n,\Rot\ov u_i(M)], 


$$ 


$\ov n$ --- нормаль к поверхности $S$, внешняя к области $V_1$; $\ov 


u{}^{\pm}_i(M)$, $N^\pm_{ni}[\ov u_i(M)]$, $i=\ov{0,1}$, ---  


граничные значения векторных функций $\ov u_i(M)$, $N_{ni}[\ov u_i(M)]$ при 


подходе к  


поверхности $S$ со стороны области $V_1$ (знак ``$+$'') или $V_0$ 


(знак ``$-$''), $[\dotsc,\dotsc]$ --- символ векторного произведения. 


 


В точках особых линий или конических точках реализация уравнений Ламе  


понимается в смысле выполнения равенства  


\begin{equation} 


\lim_{\Delta V\to M^0}\iiint\limits_{\Delta V}((\lambda_i+2\mu_i) 


\Grad\Div\ov u_i(x_1,x_2,x_3)-\mu_i\Rot\Rot\ov u_i(x_1,x_2,x_3))dv=0, 


\end{equation} 


что физически означает выполнение условий равновесия среды в особых  


точках, т.\,е. при стягивании области $\Delta V$ в особую точку $M^0$. 


 


Функция $\ov u_0(M)$, определенная в области $V_0$, при $M\to\infty$ 


принимает заданное  


значение $\ov u{}^{(0)}_0(M)$, удовлетворяющее уравнению Ламе. 


 


К такой задаче приводится ряд задач механики сплошной среды (теории  


упругости, термоупругости) [11], [13]. 


\medskip 


 


{\bf Асимптотика решений уравнений Ламе.} Изучим поведение решений  


уравнения Ламе (1) вблизи особенностей поверхностей раздела областей. 


Пусть гладкая замкнутая особая линия $L=S_1\cap S_2$, где $S_g$, $g=1,2$, 


определяется  


уравнением $f_g(x_1,x_2,x_3)=0$. Натянем на $L$ гладкую поверхность $S_0$,  


описываемую уравнением $f_0(x_1,x_2,x_3)=0$, на которой введем криволинейные  


ортогональные координаты $u$, $v$ так, что при $v=v_0=\const$ определяется 


особая  


линия $L$. Параметризуем кривую $L$ и положим $u=s$, где $s$ --- длина дуги,  


отсчитываемая от некоторой начальной точки [14]. 


Величины углов, образованных поверхностями в точках особой линии, таковы:  


\begin{align*} 


\theta_1(s)&=\arccos((\Grad f_1,\Grad f_0)/(|\Grad f_1|\,|\Grad f_0|)),\\


\theta_2(s)&=\arccos((\Grad f_2,\Grad f_0)/(|\Grad f_2|\,|\Grad f_0|)), 


\end{align*}   


здесь $(\dotsc,\dotsc)$ --- символ скалярного произведения. 


 


Пусть $\ov n\,\ov n_1\ov n_2$ --- сопровождающий трехгранник поверхности 


$S_0$ в точке $M_0\in L$,  


где $\ov n_1$, $\ov n_2$ --- орты, лежащие в плоскости, касательной к $S_0$ в 


точке $M_0$, 


$\ov n_1$ --- касательный вектор к кривой $L$, $\ov n\perp\ov n_1$, $\ov 


n\perp\ov n_2$. Точка $M$ плоскости векторов  


$\ov n$, $\ov n_2$ определяется полярными координатами $\rho=|MM_0|$, 


$\theta=\angle(\ov{M_0 M},\ov n_2)$. 


Соотношение  


\begin{equation} 


\ov r=\ov r_0+\rho(\cos\theta\ov n_2(s)+\sin\theta\ov n(s)), 


\end{equation} 


где $\ov r_0$, $\ov r$ --- радиус-векторы точек $M_0$, $M$, определяет 


переменные $\rho$, $\theta$, $s$ 


как криволинейные ортогональные координаты точки $M$ в локальной области  


особой линии с коэффициентами Ламе $h_1=1$, $h_2=\rho$, $h_3=1+\rho H_0$, 


$H_0=-(p_0\cos\theta+r_{11}\sin\theta)$, $p_0=-|\ov r_s|_v/|\ov r_v|$, 


$r_{11}=(\ov r,\ov r_{ss})/|\ov r_s|$. 


 


 


\begin{lem} 


Если характеристические уравнения 


\begin{equation} 


\sin m_q\omega=\pm m_q\sin\omega,\quad 


\mu^*\sin m_q\omega=\pm m_q\sin\omega, 


\end{equation} 


где $\mu^*=\frac{\varkappa_1\gamma-\varkappa_0}{1-\gamma}$, 


$\varkappa_i=3-4\nu_i$, $i=\ov{0,1}$, $\omega(s)=\theta_1(s)-\theta_2(s)$, 


$\gamma=\mu_0/\mu_1$, в точках особой  


линии~$L$ поверхности раздела $S$ имеют корни $m_q\in(0,1)$ и величина  


$m_5=\pi/(2\pi-\omega(s))<1$, то решение уравнения Ламе в окрестности особой 


линии  имеет асимптотику 


\begin{gather} 


\ov u=\{u_\rho,u_\theta,u_s),\quad u_\rho=\rho^{m_q}A_q+ 


o(\rho^{m_q}),\quad u_\theta=\rho^{m_q}B_q+o(\rho^{m_q}),\quad 


u_s=\rho^{m_5}V+o(\rho^{m_5}),\notag \\ 


% 


N_n[\ov u]=\{\tau_{\rho\theta},\sigma_\theta,\tau_{\theta s}\},\quad 


\tau_{\rho\theta}=\mu\rho^{m_q-1}\bigg[(m_q-1)B_q+ 


\frac{\partial A_q}{\partial\theta}\bigg]+O(1),\notag \\ 


% 


\sigma_\theta=2\mu\rho^{m_q-1}\bigg\{[\beta(1+m_q)+1]A_q+(\beta+1) 


\frac{\partial B_q}{\partial\theta}\bigg\}+O(1),\quad \tau_{\theta s}= 


\mu\rho^{m_5-1}\frac{\partial C}{\partial\theta}+O(1), 


\end{gather} 


где  


\begin{gather} 


A_q=(m_q-\varkappa)a_{1q}(s)\sin(m_q-1)\theta+(\varkappa-m_q) 


b_{1q}(s)\cos(m_q-1)\theta+\notag \\ +c_{1q}(s)\sin(m_q+1)\theta- 


d_{1q}(s)\cos(m_q+1)\theta,\notag \\ 


% 


B_q=(m_q+\varkappa)a_{1q}(s)\cos(m_q-1)\theta+(m_q+\varkappa) 


b_{1q}(s)\sin(m_q-1)\theta+\notag \\ +c_{1q}(s)\cos(m_q+1)\theta+ 


d_{1q}(s)\sin(m_q+1)\theta,\notag \\ 


% 


C=g_1(s)\cos m_5\theta+h_1(s)\sin m_5\theta,\quad 


\beta=\frac{\nu}{1-2\nu}\theta, 


\end{gather} 


$\rho$, $\theta$, $s$ --- локальные криволинейные координаты $(5)$.  


\end{lem} 


 


Здесь и далее при доказательстве леммы 1, а также леммы 2  


идентифицирующий индекс~$i$, соотносящий величины с соответствующими  


областями, с целью упрощения записи опущен, кроме величин, определяющих  


$\mu^*$ и $\gamma$. 


 


\begin{pf} 


Уравнение Ламе (1) при линейных граничных условиях (3), (4)  


допускает преобразование подобия 


\begin{equation} 


\ov u=A\ov u(Bx_1,Bx_2,Bx_3), 


\end{equation} 


где $\ov u=\{u_1,u_2,u_3\}$, $A=A(B)$, $A,B\in R$. 


Дифференцируя соотношение (9) по $B$, получим систему уравнений  


\begin{equation} 


u_\alpha\frac{\partial A}{\partial B}+A(\Grad u_\alpha,\ov r)=0,\quad 


\alpha=\ov{1,3}. 


\end{equation} 


 


Решением уравнений (10) являются однородные функции переменных  


$x_1$, $x_2$, $x_3$ \ $m$-го порядка  


\begin{equation} 


u_\alpha=F_\alpha(x_1,x_2,x_3). 


\end{equation} 


Подставляя (5) в (11), получаем представление компонент вектора $\ov u$: 


\begin{equation} 


u_\rho=\rho^{m_{\rho q}}A_q(\theta,s),\quad 


u_\theta=\rho^{m_{\theta q}}B_q(\theta,s),\quad 


u_s=\rho^{m_5}C(\theta,s), 


\end{equation} 


где $m_{\rho q}=m_{\rho q}(s)$, $m_{\theta q}=m_{\theta q}(s)$, $m_5=m_5(s)$. 


Удовлетворяя с помощью представлений (12) уравнению (4), получим  


систему дифференциальных уравнений, которая совместима при $m_{\rho 


q}=m_{\theta q}=m_q$ и  


расщепляется на систему 


\begin{equation} 


\begin{split} 


(m_q-\varkappa)\frac{\partial B_q}{\partial\theta}+(1-2\nu) 


\frac{\partial^2A_q}{\partial\theta^2}+2(1-\nu)(m^2_q-1)A_q&=0,\\ 


% 


(m_q+\varkappa)\frac{\partial A_q}{\partial\theta}+2(1-\nu) 


\frac{\partial^2B_q}{\partial\theta^2}+(1-2\nu)(m^2_q-1)B_q&=0 


\end{split} 


\end{equation} 


и уравнение 


\begin{equation} 


m^2_3C+\frac{\partial^2C}{\partial\theta^2}=0. 


\end{equation} 


Решение системы (13) и уравнения (14) дается формулами (8), откуда следует  


представление вектора $N_n[\ov u]$ в виде (7). Удовлетворяя с помощью  


представлений (8), (12) граничным условиям (3) и полагая $\theta=\theta_1(s)$ 


и $\theta=\theta_2(s)$, 


получаем однородную систему линейных алгебраических уравнений  


относительно величин $a_{1q}(s)$, $b_{1q}(s)$, $c_{1q}(s)$, $d_{1q}(s)$, а 


также однородную систему  


линейных алгебраических уравнений относительно величин $g_1(s)$, $h_1(s)$.  


Нетривиальные решения таких систем существуют при равенстве их главных  


определителей нулю, что и дает характеристические уравнения (6) и величину  


$m_5=\pi/(2\pi-\omega)$. Число положительных корней $m_q\in(0,1)$ уравнений 


(6), меньших единицы, не превышает четырех [15]. 


 


В случае плоской особой линии утверждение леммы совпадает с  


результатами работы [16]. 


\end{pf} 


 


Пусть точка $O$ является вершиной конической поверхности с прямолинейными  


образующими и гладкой криволинейной направляющей $L_0$, очерченной  


ортом $\ov r_0(s)$ образующей с началом в точке $O$. Следуя работе [17], 


введем  


локальные криволинейные координаты $\rho_1$, $\theta_1$, $s_1$ согласно 


равенству 


\begin{equation} 


\ov r_1=\rho_1(\ov r_0(s_1)\sin\theta_1+\ov n_3(s_1)\cos\theta_1), 


\end{equation} 


где $\ov r_1$ --- радиус-вектор точки $M(x_1,x_2,x_3)$, $\ov n_3(s_1)= \big[\ov 


r_0(s_1),\frac{d\ov r_0(s_1)}{ds_1}\big]$ --- нормаль к  


образующей $\ov r_0(s_1)$ конической поверхности в точке $O$, $\rho_1=|\ov 


r_1|$, $\theta_1$ --- угол,  


образованный вектором $\ov r_1$ с нормалью $\ov n_3(s_1)$. Такие координаты 


ортогональны  


с коэффициентами Ламе $h_1=1$, $h_2=\rho_1$, $h_3=\rho_1H_{0k}$, 


$H_{0k}=\sin\theta_1+d\cos\theta_1$, $d=\big|\big[\ov r_0, \frac{d^2\ov 


r_0}{ds^2_1}\big]\big|$ и  


при $\theta_1=\pi/2$ соотношения (15) описывают коническую поверхность. 


Переходя в соотношениях (11) к координатам $\rho_1$, $\theta_1$, $s_1$ 


согласно (15), получим 


\begin{equation} 


u_{\rho_1}=\rho^{m_{k1}}_1A_{0k}(\theta_1,s_1), 


\ \ u_{\theta_1}=\rho^{m_{k2}}B_{0k}(\theta_1,s_1), 


\ \ u_{s_1}=\rho^{m_{k3}}_1C_{k}(\theta_1,s_1), 


\ \ m_{kq}=m_{kq}(\theta_1,s_1),\ \ q=\ov{1,3}. 


\end{equation} 


Подставляя представления (16) в уравнение (4), записанное в координатах  


$\rho_1$, $\theta_1$, $s_1$, получим систему дифференциальных уравнений 


относительно  


неизвестных $A_{0k}(\theta_1,s_1)$, $B_{0k}(\theta_1,s_1)$, 


$C_{k}(\theta_1,s_1)$, $m_{kq}(\theta_1,s_1)$, которая из-за громоздкости  


не выписана. Для ее совместимости необходимо, чтобы 


$m_{kq}(\theta_1,s_1)=\const$. 


Рассмотрим сначала случай, когда величина $H_{0k}$ в выражении коэффициента  


Ламе $h_3$ не зависит от переменной $s_1$, что отвечает, например, круговым  


коническим поверхностям. Исходная система дифференциальных уравнений  


совместима, если $m_{k1}=m_{k2}=m_k$, и распадается на две системы: 


\begin{gather*} 


aA_{0k}+\frac{1}{H_{0k}}\frac{\partial }{\partial \theta_1} 


\bigg(H_{0k}\frac{\partial A_{0k}}{\partial \theta_1} \bigg)+ 


\frac{1}{H_{0k}}\frac{\partial }{\partial s_1} 


\bigg(\frac{\partial A_{0k}}{\partial s_1} \bigg)+b\frac{1}{H_{0k}} 


\frac{\partial }{\partial \theta_1}(H_{0k}B_{0k})=0, \\ 


% 


c\frac{\partial A_{0k}}{\partial \theta_1}+\mu_0 


\frac{\partial }{\partial \theta_1} 


\bigg(\frac{1}{H_{0k}} 


\frac{\partial }{\partial \theta_1}(H_{0k}B_{0k}) \bigg)+m_k(m_k+1) 


B_{0k}+\frac{1}{H_{0k}}\frac{\partial }{\partial s_1} 


\bigg(H_{0k}\frac{\partial B_{0k}}{\partial s_1} \bigg)=0,\\ 


% 


c\frac{1}{H_{0k}}\frac{\partial A_{0k}}{\partial s_1}+ 


\frac{\mu_0}{H_{0k}}\frac{\partial }{\partial s_1} 


\bigg(\frac{1}{H_{0k}}\frac{\partial }{\partial \theta_1} 


(H_{0k}B_{0k})\bigg)-\frac{\partial }{\partial \theta_1} 


\bigg(\frac{1}{H_{0k}}\frac{\partial B_{0k}}{\partial s_1} \bigg)=0,\\ 


% 


\frac{\partial C_k}{\partial s_1}=0,\quad \mu_0 


\frac{\partial }{\partial \theta_1} 


\bigg(\frac{1}{H_{0k}}\frac{\partial C_k}{\partial s_1} \bigg) 


-\frac{1}{H_{0k}}\frac{\partial }{\partial s_1} 


\bigg(\frac{1}{H_{0k}}\frac{\partial }{\partial \theta_1} 


(H_{0k}c_k)\bigg)=0,\\ 


% 


\frac{\mu_0}{H_{0k}}\frac{\partial }{\partial s_1} 


\bigg(\frac{1}{H_{0k}}\frac{\partial C_k}{\partial s_1} \bigg) 


+\frac{\partial }{\partial \theta_1} 


\bigg(\frac{1}{H_{0k}}\frac{\partial }{\partial \theta_1} 


(H_{0k}C_k)\bigg)+m_{k3}(m_{k3}+1)C_k=0,\\ 


% 


a=\mu_0(m_k+2)(m_k-1),\quad b=\mu_0(m_k-1)-(1+m_k),\quad 


c=\mu_0(2+m_k)-m_k,\\ 


% 


\mu_0=2(1-\nu)/(1-2\nu),\quad A_{0k}=A_{0k}(\theta_1,s_1),\quad 


B_{0k}=B_{0k}(\theta_1,s_1),\quad C_k=C_k(\theta_1,s_1).


\end{gather*} 


 


Решения таких систем получены в [17], и на основе их с учетом  


соотношений (16) находим следующие представления: 


\begin{align} 


u_{\rho_1}&=\rho^{m_k}_1(N_1P_{m_k-1}(\cos\theta_{11})+ 


N_2P_{m_k+1}(\cos\theta_{11})),\notag\\ 


% 


u_{\theta_1}&=\rho^{m_k}_1d(d_1N_1P_{m_k-1}(\cos\theta_{11})+ 


d_2N_2P_{m_k+1}(\cos\theta_{11}))/d\theta_{11},\\ 


% 


u_{s_1}&=\rho^{m_{3k}}_1N_3P^1_{m_{k3}}(\cos\theta_{11}), \notag


\end{align} 


где $P_{m_k\pm1}(\cos\theta_{11})$, $P^1_{m_{k3}}(\cos\theta_{11})$ --- 


функции Лежандра, 


$\theta_{11}=\theta_1+\alpha$, $\tg\alpha=d$, 


$H_{0k}=\sqrt{1+d^2}\sin\theta_{11}$, $N_r$, $r=\ov{1,3}$, --- произвольные 


постоянные. 


Пусть $H_{0k}$ являeтся функцией переменных $\theta_1$ и $s_1$, тогда 


исходная система  


дифференциальных уравнений совместима при $m_{k1}=m_{k2}=m_{k3}=m_k$ и 


принимает  


вид 


\begin{gather} 


aA_{0k}+\frac{1}{H_{0k}}\frac{\partial }{\partial \theta_1} 


\bigg(H_{0k}\frac{\partial A_{0k}}{\partial \theta_1} \bigg) 


+\frac{1}{H_{0k}}\frac{\partial }{\partial s_1} 


\bigg(\frac{1}{H_{0k}}\frac{\partial A_{0k}}{\partial s_1} \bigg) 


+bK=0,\notag \\ 


% 


c\frac{\partial A_{0k}}{\partial \theta_1}+m_k 


\frac{\partial K}{\partial \theta_1}+m_k(m_k+1)B_{0k} 


+\frac{1}{H_{0k}}\frac{\partial }{\partial s_1} 


\bigg(\frac{1}{H_{0k}}Q \bigg)=0,\notag\\ 


% 


c\frac{1}{H_{0k}}\frac{\partial A_{0k}}{\partial s_1} 


+\frac{m_k}{H_{0k}}\frac{\partial K}{\partial s_1} 


-\frac{\partial }{\partial \theta_1} 


\bigg(\frac{1}{H_{0k}}Q \bigg)+m_k(m_k+1)C_k=0,\\ 


% 


K=\frac{1}{H_{k0}}\bigg(\frac{\partial }{\partial \theta_1} 


(H_{k0})+\frac{\partial C_k}{\partial s_1}\bigg),\quad 


Q=\frac{\partial B_{0k}}{\partial s_1}- 


\frac{\partial }{\partial \theta_1}(H_{0k}C_k).\notag 


\end{gather} 


 


Система (18) имеет решение, приведенное в [17], и тогда  


соотношения (16) с учетом этого решения принимают вид 


\begin{gather} 


u_{\rho_1}=\rho^{m_k}_1\sum^N_{n=0}(A_{n1k}(\theta_1) 


\cos nls_1+A_{n2k}(\theta_1)\sin nls_1)+O(N^{-1}),\notag\\ 


% 


u_{\theta_1}=\rho^{m_k}_1\sum^N_{n=0}(B_{n1k}(\theta_1)\cos nls_1 


+B_{n2k}(\theta_1)\sin nls_1)+O(N^{-1}),\notag \\ 


% 


u_{s_1}=\rho^{m_k}_1\sum^N_{n=0}(C_{n1k}(\theta_1) 


\cos nls_1+C_{n2k}(\theta_1)\sin nls_1)+O(N^{-1}), \\ 


% 


A_{ngk}(\theta_1)=\sum^{2(2N+1)}_{j=1}N_{ngk} 


\Phi^{(k)}_{j,k_1}(\theta_1),\ \ 


B_{ngk}(\theta_1)=\sum^{3(2N+1)}_{j=1}N_{ngk}E_{jgk}(\theta_1),\notag \\  


C_{ngk}(\theta_1)\sum^{3(2N+1)}_{j=1}N_{ngk}F_{jgk}(\theta_1), 


\ \ g=\ov{1,2}, \notag


\end{gather} 


где $\Phi^{\{k\}}_{j,k_1}(\theta_1)$, $k_1=m_k+1;m_k-1$, --- функции 


фундаментальной системы решений,  


ограниченные при $\theta_1\in[-\pi,\pi]$; $E_{jgk}(\theta_1)$, 


$F_{jgk}(\theta_1)$ --- известные функции; $N_{ngk}$ --- 


произвольные постоянные; $l=2\pi/s_0$; $s_0$ --- длина направляющей. 


 


Подставив представления (17) или (19) в условия (4) и положив 


$\theta_1=\pi/2$,  


сгруппировав и приравняв выражения при степенях $\rho_1$, а для соотношений 


(19) и  


при гармониках $n=\ov{0,2N+1}$, нулю, получим однородные системы линейных  


алгебраических уравнений относительно произвольных постоянных.  


 


Условиями их разрешимости являются равенства их главных определителей  


нулю, что дает характеристические уравнения для показателей $m_k$ в  


представлениях (17) или (19) 


\begin{equation} 


\Delta(m_k)=0. 


\end{equation} 


Явный вид таких уравнений из-за громоздкости не выписан, а их численный  


анализ приведен в [17], [18], где установлено, что множество корней  


$m_k\in(0,1)$ уравнений вида (20) непустое. Для таких значений $m_k$ находим  


нетривиальные решения систем алгебраических уравнений.  


На основе этого устанавливаем, что компоненты вектора $N_n[\ov u]$ будут 


иметь  


особенности порядка $1-m_k$ в конической точке, и находим их явный вид. Таким  


образом, доказана 


 


 


\begin{lem} 


Если характеристическое уравнение типа $(20)$ имеет корни  


$m_k\in(0,1)$, то решение уравнения Ламе имеет асимптотику $(17)$ или $(19)$ 


$\ov u=\{u_{\rho_1},u_{\theta_1},u_{s_1}\}$ и  


\begin{gather} 


N_n[\ov u]=\{\tau_{\rho_1\theta_1},\sigma_{\theta_1}, 


\tau_{\theta_1s_1}\}, \ \ \tau_{\rho_1\theta_1}=\rho^{m_k-1}_1\mu 


\bigg[(m_k-1)B_{0k}+\frac{\partial A_{0k}}{\partial \theta_1}\bigg] 


+O(1),\notag \\ 


% 


\sigma_{\theta_1}=2\mu\rho^{m_k-1}_1\bigg\{[\beta(m_k+2)+1]A_{0k} 


+(\beta+1)\frac{\partial B_{0k}}{\partial \theta_1}+\notag \\ 


+\beta H^{-1}_{0k}\frac{\partial C_k}{\partial s_1} 


+\beta B_{0k}H^{-1}_{0k} 


\frac{\partial H_{0k}}{\partial \theta_1}\bigg\}+ 


O(1),\ \ \beta=\nu/(1-2\nu),\notag\\ 


% 


\tau_{\theta_1s_1}=\mu\rho^{m_k-1}_1\bigg(-H^{-1}_{0k} 


\frac{\partial H_{0k}}{\partial \theta_1}C_k+ 


\frac{\partial C_k}{\partial \theta_1}+H^{-1}_{0k} 


\frac{\partial B_{0k}}{\partial s_1}\bigg)+O(1).


\end{gather} 


\end{lem} 


 


Отметим, что здесь, как и в лемме 1, идентифицирующий индекс $i$,  


соотносящий величины с соответствующей областью, для упрощения записи  


опущен. Компоненты векторов $\ov u$, $N_n[\ov u]$ зависят от $\mu_0$, 


$\mu_1$, а величина $m_k$ зависит  


от отношения $\gamma=\mu_0/\mu_1$ посредством соответствующего 


характеристического уравнения, что детально показано в работе [17]. 


\medskip 


 


{\bf Разрешимость задачи.} Решение уравнения Ламе (1) представляется  


обобщенными упругими потенциалами простого и двойного слоя [10].  


 


\begin{lem} 


Если в точках гладкости замкнутой кусочно-гладкой поверхности  


$S$ плотность $\ov\varphi_2(M)$ обобщенного упругого потенциала двойного слоя  


$\ov W(M)$ удовлетворяет условию Липшица--Г\"ельдера, а в особых точках $M^0$ 


имеет место представление 


\begin{equation} 


\ov\varphi_2(M)=O(|MM^0|^\alpha),\quad \alpha\in(0,1),\quad 


\ov\varphi_2(M^0)=0, 


\end{equation} 


то обобщенный упругий потенциал двойного слоя существует и в точках  


гладкости поверхности имеет место обычное представление 


\begin{equation} 


\ov W{}^\pm(M_0)=\mp\ov\varphi_2(M_0)+\ov W(M_0), 


\end{equation} 


где $\ov W{}^\pm(M_0)$ --- граничное значение потенциала при стремлении точки  


$M\in V_0$ $($знак ``$-$''$)$ или точки $M\in V_1$ $($знак ``$+$''$)$ к точке 


$M_0\in S;$ 


$\ov W(M_0)$ --- прямое значение потенциала в $M_0$. 


\end{lem} 


 


\begin{pf} 


Пусть поверхность $S$ содержит одну особую точку $O_1$,  


являющуюся, например, конической точкой, а точка $M_0\in S$ является точкой  


гладкости. Доказательство в случае множества особых точек аналогично.  


Построим сферу с центром в особой точке $O_1$ радиуса $R_0<|M_0O_1|$. Тогда  


$S=S_1\cup S_2$, где $S_1$ --- часть поверхности $S$, лежащей вне сферы и 


являющейся  


гладкой разомкнутой поверхностью, а $S_2$ --- часть поверхности $S$, 


содержащаяся внутри такой сферы. Потенциал представляется в виде суммы 


\begin{equation} 


\ov W(M)=\iint\limits_S \Gamma_2(M,N)\ov\varphi_2(N)ds_N= 


\iint\limits_{S_1}\Gamma_2(M,N)\ov\varphi_2(N)ds_N+ 


\iint\limits_{S_2}\Gamma_2(M,N)\ov\varphi_2(N)ds_N, 


\end{equation} 


где 


\begin{gather*} 


\Gamma_2(M,N)=\|\Gamma_{2mp}(M,N)\|,\quad m,p=\ov{1,3};\\ 


% 


\Gamma_{2mp}(M,N)=\bigg[m_{01}\delta_{mp}+ 


\frac{n_{01}(y_m-x_m)(y_p-x_p)}{r^2}\bigg] 


\sum^3_{i=1}\frac{(x_l-y_l)n_l(N)}{r^3}+ \\ 


% 


+m_{01}\bigg[n_m(N)\frac{x_p-y_p}{r^3}-n_p(N) 


\frac{x_m-y_m}{r^3}\bigg],\quad m=\ov{1,3};\ \ p=\ov{1,3},\\ 


% 


m_{01}=\frac{1}{2\pi}\frac{\mu}{\lambda+2\mu},\ \  


n_{01}=\frac{3}{2\pi}\frac{\lambda+\mu}{\lambda+2\mu},\ \  


\lambda=\frac{E\nu}{(1+\nu)(1-2\nu)},\ \  


\mu=\frac{E}{2(1+\nu)},\ \ N=N(y_1,y_2,y_3), 


\end{gather*} 


$r=((x_1-y_1)^2+(x_2-y_2)^2+(x_3-y_3)^2)^{1/2}$ --- матрица потенциала 


двойного слоя, $\delta_{mp}$ --- символ Кронекера. 


 


Первое слагаемое в правой части (24) существует как обобщенный упругий  


потенциал двойного слоя на разомкнутой поверхности Ляпунова [10]. Второе  


слагаемое содержит интегралы, которые при переходе к переменным $\rho_1$, 


$s_1$ (15) вычисляются как несобственные интегралы типа 


$$ 


\int^{R_0}_0d\rho_1\int^{s_1}_0\frac{A_{s_1}(s_1)}{\rho^{1-\alpha}_1} 


ds_1=\frac{R^\alpha_0}{\alpha}\int^{s_1}_0 A_{s_1}(s_1)ds_1<\infty. 


$$ 


 


Пусть точка $M$ стремится к точке $M_0\in S_1$, тогда, учитывая для первого  


слагаемого в представлении (24) известные граничные значения обобщенного  


упругого потенциала двойного слоя на разомкнутых поверхностях Ляпунова  


([13], с.\,218) и непрерывность второго слагаемого при $M_0\in S_1$, получаем 


формулу (23). 


\end{pf} 


 


\begin{lem} 


Если в точках гладкости кусочно-гладкой поверхности $S$ 


плотность $\ov\varphi_1(M)$ обобщенного упругого потенциала простого слоя  


$\ov V(M)$


удовлетворяет условию Липшица--Г\"ельдера, а в особых точках $M^0$ стремится 


к бесконечности и имеет место асимптотическое представление  


\begin{equation} 


\ov\varphi_1(M)=O\bigg(\frac{1}{|MM^0|^\alpha} \bigg),\quad 


\alpha\in(0,1), 


\end{equation} 


то обобщенный упругий потенциал простого слоя существует и является  


непрерывной функцией при переходе поверхности $S$ в ее точках гладкости и в  


этих точках имеет место обычное представление  


\begin{equation} 


N^\pm_n[\ov V(M_0)]=\pm\ov\varphi_1(M_0)+N_n[\ov V(M_0)], 


\end{equation} 


где $N^\pm_n[\ov V(M_0)]$ --- граничные значения вектора $N_n[\ov V(M_0)]$


при подходе точки $M$ 


к точке $M_0\in S$ со стороны области $V_0$ $($знак ``$-$''$)$ или $V_1$ 


$($знак ``$+$''$);$ $N_n[\ov V(M_0)]$ --- 


прямое значение оператора от потенциала на поверхности $S$. 


\end{lem} 


 


\begin{pf} 


Аналогично предыдущему проиллюстрируем  


доказательство на примере наличия на поверхности одной особой точки, также  


конической. Представляя, как и в предыдущем доказательстве, поверхность  


$S=S_1\cup S_2$, получаем 


\begin{equation} 


\ov V(M)=\iint\limits_S \Gamma(M,N)\ov\varphi_1(N)ds_N= 


\iint\limits_{S_1}\Gamma(M,N)\ov\varphi_1(N)ds_N+ 


\iint\limits_{S_2}\Gamma(M,N)\ov\varphi_1(N)ds_N, 


\end{equation} 


где $\Gamma(M,N)=\|\Gamma_{mp}(M,N)\|$, 


$$ 


\Gamma_{mp}(M,N)=\frac{1}{4\pi\mu(\lambda+2\mu)}\bigg[(\lambda+\mu) 


\frac{\partial r}{\partial x_m}\frac{\partial r}{\partial x_p} 


+(\lambda+3\mu)\delta_{mp}\bigg]\frac{1}{r} 


$$ 


--- матрица потенциала простого слоя. 


Первый интеграл правой части соотношения (27) существует как интеграл по  


поверхности Ляпунова. Переходя в интеграле второго слагаемого с учетом  


(25) к переменным $\rho_1$, $s_1$ с помощью (15), убеждаемся, что он 


существует как  


несобственный. Устремляя в представлении (27) $M$ к $M_0$ и учитывая  


непрерывность обобщенного упругого потенциала простого слоя при переходе  


разомкнутых поверхностей Ляпунова ([10], с.\,547), 


а также непрерывность второго  


слагаемого, убеждаемся, что $\ov V(M)$ является непрерывной функцией при  


переходе поверхности $S$ в ее точках гладкости. 


 


Подействуем оператором $N_n[\ ]$ на равенство (27), перейдем в полученном  


равенстве к пределу при $M\to M_0$. Принимая во внимание известные  


граничные значения для $N_n[\ov V(M)]$ на разомкнутой поверхности Ляпунова 


([13], с.\,231) и непрерывность второго слагаемого правой части, 


получаем утверждение леммы. При этом интеграл $N_n[\ov V(M_0)]$ понимаем как  


сингулярный. 


\end{pf} 


 


Перейдем к общему случаю разрешимости задачи.  


Пусть 


\begin{gather*} 


\Gamma_1(M,N)=\|\Gamma_{1mp}(M,N)\|,\\  


%


\Gamma_{1mp}(M,N)=-0,5\bigg[m_{01}\delta_{pm}+n_{01} 


\frac{(x_m-y_m)(x_p-y_p)}{r^2}\bigg]\frac{1}{r^3} 


\sum^3_{l=1}(x_l-y_l)n_l(M)+\\ 


% 


+m_{01}\bigg[n_m(M)\frac{x_p-y_p}{r^3}-n_p(M) 


\frac{x_m-y_m}{r^3}\bigg],\ \ \Gamma_3(M,N)=\|\Gamma_{3mp}(M,N)\|,\\  


%


\Gamma_{3mp}(M,N)=\sum^3_{s=1}(\sigma^{(p)}_{ms}n_s),\\ 


% 


\sigma^{(p)}_{ms}=2\mu\varepsilon^{(p)}_{ms}+\lambda\delta_{ms} 


\theta^{(p)},\ \ \varepsilon^{(p)}_{ms}=0,5\bigg( 


\frac{\partial u^{(p)}_s}{\partial x_m}+ 


\frac{\partial u^{(p)}_m}{\partial x_s}\bigg),\\  


%


u^{(h)}_g=\Gamma_{2gh},\ \ g,h=\ov{1,3};\ \ p=\ov{1,3},\ \  


\theta^{(p)}=\varepsilon^{(p)}_{11}+ 


\varepsilon^{(p)}_{22}+\varepsilon^{(p)}_{33}, 


\end{gather*} 


матрицы $\Gamma^{(i)}(M,N)$, $\Gamma^{(i)}_t(M,N)$, $t=\ov{1,3}$; 


$i=\ov{0,1}$, следуют  


соответственно из матриц $\Gamma(M,N)$, $\Gamma_t(M,N)$, $t=\ov{1,3}$, если в 


элементах последних  


заменить постоянные $\nu$, $\mu$ соответственно на $\nu_i$, $\mu_i$, т.\,е. 


при значении  


индекса $i=0$ величины отвечают области $V_0$, а при $i=1$ --- области $V_1$, 


$M_0$ --- точки гладкости поверхности $S$. 


 


 


\begin{thmnonum} 


Если в особых точках кусочно-гладкой поверхности корни  


характеристических уравнений $(6)$, $(20)$ принадлежат интервалу $(0,1)$ 


и имеют место следующие соотношения\/${:}$ 


\begin{align} %% return ali by =


&0,5\sum^3_{p=1}\iint\limits_S\{|\Gamma^{(1)}_{mp}(M,N)- 


\Gamma^{(0)}_{mp}(M,N)|+|\Gamma^{(1)}_{2mp}(M,N)- 


\Gamma^{(0)}_{2mp}(M,N)|\}ds_n=A_m<1,\\ 


% 


&0,5\sum^3_{p=1}\iint\limits_S\{|\Gamma^{(1)}_{1mp}(M,N)- 


\Gamma^{(0)}_{1mp}(M,N)|+|\Gamma^{(1)}_{3mp}(M,N)- 


\Gamma^{(0)}_{3mp}(M,N)|\}ds_n=B_m<1,\ \ m=\ov{1,3},\notag


\end{align} 


то задача сопряжения разрешима. 


\end{thmnonum} 


 


\begin{pf} 


Пусть конечная односвязная область ограничена  


поверхностью $S=\mathop{\cup}\limits^{n_1}_{j=1}S_j$, которая содержит 


конические точки $O_l$, $l=\ov{1,m_1}$, и  


гладкие замкнутые непересекающиеся особые линии $L_j=S_j\cap S_{j+1}$, 


$j=\ov{1,n_1-1}$, $O_l\notin L_j$.  


Поверхности $S_j$ задаются уравнениями $f_j(x_1,x_2,x_3)=0$, а натянутые на  


контуры $L_j$, определяются уравнениями $f_{0j}(x_1,x_2,x_3)=0$. 


 


Учитывая характер решений, установленный в леммах 1, 2, представляем  


решения уравнений Ламе в виде 


\begin{equation} 


\ov u_0(M)=\ov u_{01}(M)+\ov u_{02}(M),\quad 


\ov u_1(M)=\ov u_{11}(M)+\ov u_{12}(M), 


\end{equation} 


где $\ov u_{i1}(M)$, $i=\ov{0,1}$, реализуют асимптотику, приведенную в 


леммах 1 и 2, а  


$\ov u_{i2}(M)$ являются ограниченными непрерывными функциями в своих 


областях  


определения. 


 


Согласно леммам 3 и 4 характер обобщенного упругого потенциала двойного  


слоя и оператора $N_n[\ ]$ от обобщенного упругого потенциала простого слоя  


вблизи точек гладкости поверхности $S$ определяется плотностями потенциалов  


с точностью до величин прямых значений их на поверхности. 


 


Представим сингулярные составляющие посредством обобщенных упругих  


потенциалов простого и двойного слоя 


\begin{equation} 


\ov u_{i1}(M)=\iint\limits_S \Gamma^{(i)}(M,N)\ov\varphi_{i11}(N)ds_N+ 


\iint\limits_S \Gamma^{(i)}_2(M,N)\ov\varphi_{i12}(N)ds_N, 


\end{equation} 


где значение индекса $i=0$ относится к величинам области $V_0=R_3\setminus 


V_1$, а $i=1$ --- 


области $V_1$. Векторы плотностей представляем в виде 


\begin{gather} 


\ov\varphi_{i11}(N)=\{\varphi_{i111}(N),\varphi_{i112}(N), 


\varphi_{i113}(N)\},\quad \ov\varphi_{i12}(N)=\{\varphi_{i121}(N), 


\varphi_{i122}(N),\varphi_{i123}(N)\},\\ 


% 


\varphi_{i111}(N)=\bigg(\prod^{n_1-1}_{j=1}\varphi^{(1)}_{i11j}(N)\bigg) 


\bigg(\prod^{m_1}_{k=1}\varphi^{(2)}_{i11k}(N)\bigg),\quad 


\varphi_{i121}(N)\bigg(\prod^{n_1-1}_{j=1} 


\varphi^{(3)}_{i21j}(N)\bigg)\bigg(\prod^{m_1}_{k=1} 


\varphi^{(4)}_{i21k}(N)\bigg),\notag \\ 


% 


M=M(x_1,x_2,x_3),\quad N=N(y_1,y_2,y_3).\notag 


\end{gather} 


 


В формулах (31) функции $\varphi^{(1)}_{i11j}(N)$, $\varphi^{(3)}_{i21j}(N)$ 


определяются геометрией  


поверхностей вблизи особых линий и согласно леммам 1, 3 ,4 имеют вид 


\begin{gather} 


\varphi^{(1)}_{i11j}(N)=2\mu_i(-1)^{i+1}\sum^{p_2}_{q=1}\Big\langle 


R^{m_{qj}-1}_j(N)\bigg\{(m_{qj}+1)B_{ijq}[T_j(N)]+ 


\frac{\partial A_{ijq}[T_j(N)]}{\partial T_j(N)}\bigg\}\Big\rangle,\notag\\ 


% 


\varphi^{(1)}_{i12j}(N)=2\mu_i(-1)^{i+1}\sum^{p_2}_{q=1}\Big\langle 


R^{m_{q1j}-1}_j(N)\bigg\{[\beta_i(1+n_i)+1]A_{ijq}[T_j(N)]+(\beta_i+1) 


\frac{\partial B_{ijq}[T_j(N)]}{\partial T_j(N)}\bigg\}\Big\rangle,\notag\\ 


% 


\displaybreak[3]


\varphi^{(1)}_{i13j}(N)=(-1)^i \mu_i R^{m_{5j}-1}_j(N) 


\frac{\partial C_{ij}[T_j(N)]}{\partial T_j(N)},\quad 


\varphi^{(3)}_{i12j}(N)=(-1)^i\sum^{p_2}_{q=1}R^{m_{qj}}_j 


(N)A_{ijq}[T_j(N)],\notag\\ 


% 


\varphi^{(3)}_{i22j}(N)=(-1)^i\sum^{p_2}_{q=1}R^{m_{qj}}_j(N)B_{ijq} 


[T_j(N)],\quad \varphi^{(3)}_{i23j}(N)=(-1)^i R^{m_{3j}}_j 


(N)C_{ij}[T_j(N)],\\ 


% 


A_{ijq}(T_j(N))=(m_{qj}-\varkappa_i)a_{1ijq}\sin((m_{qj}-1) 


T_j(N))+(\varkappa_i-m_{qj})b_{1ijq}\cos((m_{qj}-1)T_j(N))+\notag\\ 


% 


+c_{1ijq}\sin((m_{qj}+1)T_j(N))-d_{1ijq}\cos((m_{qj}+1)T_j(N)),\notag\\ 


% 


B_{ijq}(T_j(N))=(m_{qj}+\varkappa_i)a_{1ijq}\cos((m_{qj}-1)T_j(N))+ 


(m_{qj}+\varkappa_i)b_{1ijq}\sin((m_{qj}-1)T_j(N))+\notag \\ 


% 


+c_{1ijq}\cos((m_{qj}+1)T_j(N))+d_{1ijq}\sin((m_{qj}+1)T_j(N)),\notag\\ 


% 


C_{ij}(T_j(N))=g_{1ij}\cos(m_{5j}T_j(N))+h_{1ij} 


\sin(m_{5j}T_j(N));\notag\\ 


% 


R_j(N)=((f_{0j}(N)/|\Grad f_{0j}(N)|)^2+Q^2_j(N))^{1/2},\quad 


T_j(N)=\arcsin(f_{0j}(N)/(|\Grad f_{0j}(N)|R_j(N))),\notag \\ 


% 


Q_j(N)=(f_j(N)/|\Grad f_j(N)|)+(1-(\ov G_{2j}(N),\ov G_{0j}(N))^2 


f_{0j}(N)/|\Grad f_{0j}(N)|(\ov G_{0j}(N),\ov G_{2j}(N))),\notag\\ 


% 


\ov G_j(N)=\Grad f_j(N)/|\Grad f_j(N)|,\quad 


\ov G_{0j}(N)=\Grad f_{0j}(N)/|\Grad f_{0j}(N)|,\notag \\ 


% 


\ov G_{2j}(N)=[[\ov G_j(N),\ov G_{j+1}(N)]/ 


|[\ov G_j(N),\ov G_{j+1}(N)]|,\ov G_{0j}(N)],\notag 


\end{gather} 


$p_2$ --- число корней $m_{qj}\in(0,1)$ характеристических  


уравнений (6). 


 


Функции $\varphi^{(2)}_{i11k}(N)$, $\varphi^{(4)}_{i21k}(N)$ определяются 


геометрией поверхностей вблизи  


конических точек и согласно леммам 2--4 имеют вид (17) и (21) или (19) 


и (21), в которых необходимо заменить $\rho_1$, $\theta_1$, $s_1$ 


соответственно на 


\begin{gather*} 


R_l(N)=((y_1-x_{1l})^2+(y_2-x_{2l})^2+(y_3-x_{3l})^2)^{1/2},\\ 


% 


T_{ln}(N)=\arcsin(R^{-1}_l(N)(l_{0l}(y_1-x_{1l})+m_{0l}(y_2-x_{2l})+ 


n_{0l}(y_3-x_{3l}))),\\ 


% 


S_l(N)=\int^{\varphi_l}_0((dl_{0l}(N)/d\varphi_l)^2+ 


(dm_{0l}(N)/d\varphi_l)^2+(dn_{0l}(N)/d\varphi_l)^2)d\varphi_l, 


\end{gather*} 


где $\varphi_l=\varphi_l(y_2,y_3)$, $O_l(x_{1l},x_{2l},x_{3l})$, $l_{0k}(N)$, 


$m_{0k}(N)$, $n_{0k}(N)$ ---  


координаты орта образующей $\ov r_{0k}(N)$, и просуммировать в  


зависимости от числа корней $m_k\in(0,1)$ уравнения (20). 


 


Представление сингулярных составляющих обобщенными упругими  


потенциалами двойного и простого слоя с плотностями (31), (32) отвечает  


условиям лемм 3 и 4. Удовлетворяя с помощью (29), (30) граничным условиям  


(4) в особых точках при переходе к соответствующим локальным координатам,  


убеждаемся, что они реализуются.  


 


Регулярные составляющие решений уравнений Ламе в формулах (29)  


представляем также обобщенными упругими потенциалами простого и  


двойного слоя с неизвестными плотностями $\ov\varphi_{i21}(N)$, 


$\ov\varphi_{i22}(N)$, $i=\ov{0,1}$, 


\begin{align} 


\ov u_{02}(M)&=\ov u{}^{(0)}_0(M)+\iint\limits_S\Gamma(M,N) 


\ov\varphi_{021}(N)ds_N+\iint\limits_S\Gamma_2(M,N) 


\ov\varphi_{022}(N)ds_N,\\ 


% 


\ov u_{12}(M)&=\ov u{}^{(0)}_1(M)+\iint\limits_S\Gamma(M,N) 


\ov\varphi_{121}(N)ds_N+\iint\limits_S\Gamma_2(M,N) 


\ov\varphi_{122}(N)ds_N,


\end{align} 


где $\ov u{}^{(0)}_i(M)$ --- известные векторы, удовлетворяющие уравнению 


Ламе (1) и  


обусловленные физической сущностью задачи, а также условиям на  


бесконечности. 


 


Подставляя представления (29)--(34) в граничные условия (3) в точках  


гладкости и учитывая формулы (23), (26), получим систему интегральных  


уравнений, в которой полагаем ввиду произвольности плотностей 


$$ 


\ov\varphi_{021}(N)=\ov\varphi_{121}(N),\quad 


\ov\varphi_{022}(N)=\ov\varphi_{122}(N). 


$$ 


В результате придем к системе интегральных уравнений 


\begin{align} 


\ov\varphi_{022}(M_0)&-0,5\iint\limits_S[\Gamma^{(1)}(M_0,N)- 


\Gamma^{(0)}(M_0,N)]\ov\varphi_{021}(N)ds_N- \notag \\


% 


&\qquad-0,5\iint\limits_S[\Gamma^{(1)}_2(M_0,N)-\Gamma^{(0)}_2 


(M_0,N)]\ov\varphi_{022}(N)ds_N=-0,5\ov g_1(M_0),\notag \\ 


% 


\ov\varphi_{021}(M_0)&+0,5\iint\limits_S[\Gamma^{(1)}_1(M_0,N)- 


\Gamma^{(0)}_1(M_0,N)]\ov\varphi_{021}(N)ds_N+ \notag \\


% 


&\qquad+0,5\iint\limits_S[\Gamma^{(1)}_3(M_0,N)-\Gamma^{(0)}_3 


(M_0,N)]\ov\varphi_{022}(N)ds_N=0,5\ov h_1(M_0),\end{align} 


где 


\begin{gather*} 


\ov g_1(M_0)=\ov u{}^-_{01}(M_0)+\ov u{}^{(0)-}_0(M_0)- 


(\ov u{}^+_{11}(M_0)+\ov u{}^{(0)+}_1(M_0)), \\ 


\ov h_1(M_0)=N^-_{n0}[\ov u_{01}(M_0)+\ov u{}^{(0)}_0(M_0)]- 


N^+_{n1}[\ov u_{11}(M_0)+\ov u{}^{(0)}_1(M_0)], 


\end{gather*} 


матрицы $\Gamma^{(i)}(M,N)$, $\Gamma^{(i)}_t(M,N)$, $t=\ov{1,3}$, и оператор 


$N_{ni}[\ ]$ 


при значении индекса $i=0$ отвечают области $V_0$, а при $i=1$ --- области 


$V_1$, $M_0$ --- точки гладкости поверхности $S$. 


Ввиду выполнения условий (3) в особых точках правые части уравнений  


(35) являются непрерывными ограниченными функциями на поверхности $S$, а  


матрицы $\Gamma^{(i)}_1(M_0,N)$, $\Gamma^{(i)}_3(M_0,N)$ порождают 


сингулярные интегралы, поэтому (35)  


является системой сингулярных интегральных уравнений [19]. 


 


Решение системы (35) аналогично [13] представляем в итерационной форме 


\begin{align} 


\ov\varphi{}^{(0)}_{022}(M)&=-0,5\ov g_1(M),\quad 


\ov\varphi{}^{(0)}_{021}=0,5\ov h_1(M),\notag\\ 


% 


\ov\varphi{}^{(n)}_{022}(M)&=0,5\iint\limits_S[\Gamma^{(1)}(M,N)- 


\Gamma^{(0)}(M,N)]\ov\varphi{}^{(n-1)}_{021}(N)ds_N+\notag\\ 


%


&\qquad+0,5\iint\limits_S[\Gamma^{(1)}_2(M,N)- 


\Gamma^{(0)}_2(M,N)]\ov\varphi{}^{(n-1)}_{022}(N) 


ds_N-0,5\ov g_1(M),\notag \\ 


% 


\ov\varphi{}^{(n)}_{021}(M)&=-0,5\iint\limits_S[\Gamma^{(1)}_1(M,N)- 


\Gamma^{(0)}_1(M,N)]\ov\varphi{}^{(n-1)}_{021}(N)ds_N-\notag \\ 


% 


&\qquad-0,5\iint\limits_S[\Gamma^{(1)}_3(M,N)- 


\Gamma^{(0)}_3(M,N)]\ov\varphi{}^{(n-1)}_{022}(N)ds_N+0,5\ov h_1(M),\\ 


%


\ov\varphi_{022}(M)&=\lim_{n\to\infty}\ov\varphi{}^{(n)}_{022}(M),\quad 


\ov\varphi_{021}(M)=\lim_{n\to\infty}\ov\varphi{}^{(n)}_{021}(M). 


\end{align} 


Здесь точка $M$ принадлежит множеству точек гладкости поверхности $S$. 


Сходимость пределов (37) равносильна сходимости рядов 


\begin{equation} 


\ov\varphi{}^{(0)}_{021}(M)+\sum^\infty_{n=1}[ 


\ov\varphi{}^{(n)}_{021}(M)-\ov\varphi{}^{(n-1)}_{021}(M)],\quad 


\ov\varphi{}^{(0)}_{022}(M)+\sum^\infty_{n=1} 


[\ov\varphi{}^{(n)}_{022}(M)-\ov\varphi{}^{(n-1)}_{022}(M)]. 


\end{equation} 


 


Переходя к скалярным компонентам (38) и учитывая условия (28), находим,  


что они мажорируются сходящимися числовыми рядами 


$\sum\limits^\infty_{n=1}A^n_m M_{1m}$, $\sum\limits^\infty_{n=1}B^n_m 


M_{1m}$, где $M_{1m}=\max\{|0,5h_{1m}(M)|$, $|0,5g_{1m}(M)|\}$, $\ov 


%%%%%% join in eng.                      ^^^^


h_1(M)=\{h_{11}(M),h_{12}(M),h_{13}(M)\}$, $\ov g_1(M)=\{g_{11}(M), 


g_{12}(M),g_{13}(M)\}$, 


и таким образом, ряды (38) сходятся абсолютно и  


равномерно в точках гладкости поверхности $S$. Выполняя предельный переход  


в (36) с учетом равномерной сходимости (38) в точках гладкости $M,N\in S$,  


получаем, что (36), (37) является решением системы (35). 


 


Пусть система имеет два решения $\ov\varphi{}^{(a)}_{021}(M)$, 


$\ov\varphi{}^{(a)}_{022}(M)$ и $\ov\varphi{}^{(b)}_{021}(M)$, 


$\ov\varphi{}^{(b)}_{022}(M)$, тогда, 


подставляя их в (35) и вычитая уравнения, получим 


\begin{equation} 


\begin{split} 


&-\ov u_b(M_0)+0,5\iint\limits_S[\Gamma^{(1)}(M_0,N)- 


\Gamma^{(0)}(M_0,N)]\ov u_a(N)ds_N+ \\ &\qquad+


0,5\iint\limits_S[\Gamma^{(1)}_2(M_0,N)- 


\Gamma^{(0)}_2(M_0,N)]\ov u_b(N)ds_N=0,\\ 


% 


&-\ov u_a(M_0)+0,5\iint\limits_S[\Gamma^{(1)}_1(M_0,N)- 


\Gamma^{(0)}_1(M_0,N)]\ov u_a(N)ds_N+ \\ &\qquad +


0,5\iint\limits_S[\Gamma^{(1)}_3(M_0,N)- 


\Gamma^{(0)}_3(M_0,N)]\ov u_b(N)ds_N=0,


\end{split} 


\end{equation} 


где  


\begin{align*} 


\ov u_a(M)&=\{u_{a1}(M),u_{a2}(M),u_{a3}(M)\}= 


\ov\varphi{}^{(a)}_{021}(M)-\ov\varphi{}^{(b)}_{021}(M),\\ 


% 


\ov u_b(M)&=\{u_{b1}(M),u_{b2}(M),u_{b3}(M)\}= 


\ov\varphi{}^{(a)}_{022}(M)-\ov\varphi{}^{(b)}_{022}(M). 


\end{align*} 


Решения системы (35) ограничены на $S$ и поэтому $|u_{ql}(M)|<L$, 


$q=\ov{1,2}$, $l=\ov{1,3}$. 


Тогда из (39) с учетом (28) имеем $L\le|A_m|L$, $L\le|B_m|L$, откуда следует 


$|A_m|\ge1$, $|B_m|\ge1$,  


что противоречит условию теоремы, т.\,е. решение задачи  


единственно. 


\end{pf} 


 


Для приведения задачи сопряжения к системе сингулярных интегральных  


уравнений в случае $N$ попарно непересекающихся конечных односвязных  


областей $V_i$ сингулярные составляющие решений уравнений Ламе строим  


аналогично (30)--(32) для каждой области. Регулярные слагаемые берем, следуя  


(34) для каждой области $V_i$, $i=\ov{1,N}$, а для области 


$V_0=R_3\setminus\mathop{\cup}\limits^N_{i=1}V_i$ представляем  


в виде $\ov u_{01}(M)=\sum\limits^N_{j=1}\ov u_{0j1}(M)$, $\ov 


u_{02}(M)=\sum\limits^N_{j=1}\ov u_{0j2}(M)$, 


где $\ov u_{0j1}(M)$ строится с учетом  


негладкостей поверхностей раздела $S_i$ областей $V_i$ и $V_0$ аналогично  


изложенному выше, а слагаемые $\ov u_{0j2}(M)$ представляются в виде (33), 


(34). 


 


 


\begin{thebibliography}{99} 


 


\bibitem{1} 


Денисюк И.Т. {\em Решение одной задачи сопряжения для составной области с  


угловыми точками на линиях раздела} // Изв. вузов. Математика. -- 1996. -- 


\No\,6. -- С.\,17--24. 


\bibitem{2} 


Денисюк И.Т. {\em Одна задача сопряжения аналитических функций в  


аффинно-преобразованных областях с кусочно-гладкими границами} // Изв. вузов.  


Математика. -- 2000. -- \No\,6. -- С.\,70--74. 


\bibitem{3} 


Денисюк И.Т. {\em Термоупругость изотропной пластинки с угловыми  


включениями} // Изв. РАН. Сер. механ. твердого тела. -- 1999. -- 


\No\,2. -- С.\,148--155. 


\bibitem{4} 


Денисюк И.Т. {\em Одна модель тонких упругих включений в изотропной  


пластинке} // Изв. РАН. Сер. механ. твердого тела. -- 2000. -- 


\No\,4. -- С.\,140--148. 


\bibitem{5} 


Денисюк И.Т. {\em Напряжения анизотропной пластинки с угловыми  


включениями} // Прикл. механика. -- 1999. -- \No\,2. -- С.\,76--84. 


\bibitem{6} 


Денисюк И.Т. {\em Особенность напряжений анизотропной пластинки с угловым  


вырезом} // Прикл. механика. -- 1996. -- \No\,1. -- С.\,48--52. 


\bibitem{7} 


Денисюк I.Т. {\em Термонапруження в анiзотропнiй пластинi з кутовими  


включеннями} // Фiз.-хiм. механ. матерiалiв. -- 1999. -- \No\,3. -- 


С.\,59--68. 


\bibitem{8} 


Барабаш С.С., Денисюк I.Т. {\em Поздовжнiй зсув анiзотропного тiла з кутовими  


вкрапленнями} // Фiз.-хiм. механ. матерiалiв. -- 1996. -- \No\,4. -- 


С.\,86--90. 


\bibitem{9} 


Денисюк I.Т. {\em Поздовжнiй зсув анiзотропного тiла з пружними смугами} //  


Фiз.-хiм. механ. матерiалiв. -- 1997. -- \No\,1. -- С\,.51--56. 


\bibitem{10} 


Партон В.З., Перлин П.И. {\em Методы математической теории упругости}. -- М.:  


Наука, 1981. -- 688~с. 


\bibitem{11} 


Седов Л.И. {\em Механика сплошной среды}. T.\,1. -- М.: Наука, 1973. --  


536~с. 


\bibitem{12} 


Седов Л.И. {\em Механика сплошной среды}. T.\,2. -- М.: Наука, 1973. --  


568~с. 


\bibitem{13} 


Купрадзе В.Д., Гегелиа Т.Г., Башелейшвили М.О., Бурчуладзе Т.В.  


{\em Трехмерные задачи математической теории упругости}. -- Тбилиси: 


Изд-во Тбилисск. ун-та, 1968. -- 627~с. 


\bibitem{14} 


Норден А.П. {\em Краткий курс дифференциальной геометрии}. -- М.:  


ГИФМЛ, 1958. -- 244~с. 


\bibitem{15} 


Денисюк I.Т. {\em Сингулярнi напруження в iзотропнiй матрицi з пружним  


клином} // Фiз.--хiм. механ. матерiалiв. -- 1992. -- \No\,4. -- С.\,76--81. 


\bibitem{16} 


Денисюк И.Т. {\em Напряженное состояние вблизи особой линии поверхности  


раздела сред} // Изв. РАН. Сер. Механ. твердого тела. -- 1995. -- \No\,5. 


-- С.\,64--70. 


\bibitem{17} 


Денисюк И.Т. {\em Напряжения вблизи конической точки поверхности раздела  


сред} // Изв. РАН. Сер. механ. твердого тела. -- 2001. -- \No\,3. 


-- С.\,68--77. 


\bibitem{18} 


Денисюк I.Т. {\em Напруження бiля конiчних та пiрамiдальних включень} 


// Фiз.-хiм. механ. матерiалiв. -- 2000. -- \No\,3. -- С.\,16--20. 


\bibitem{19} 


Михлин С.Г. {\em Многомерные сингулярные интегралы и интегральные  


уравнения}. -- М.: Физматгиз, 1962. -- 254~с. 


 


\end{thebibliography} 


 


\vskip 0.5cm 


 


\postup{\it Луцкий государственный}{\qquad \it Поступили} 


\postup{\it технический университет $($Украина$)$}{\it первый вариант 


$06.12.2001$} 


\postup{}{\it окончательный вариант $08.04.2003$} 


 


\label{end} 


\end{document} 


